\section{Method}
\label{sec:method}

\subsection{Problem Formulation}

We formulate mapless navigation as a continuous-control reinforcement
learning problem. The state vector
$s = [l_1, \ldots, l_{180}, d, \cos\theta, \sin\theta, v_{prev}, \omega_{prev}]
\in \mathbb{R}^{185}$
combines a 180-beam 2D LiDAR scan (7\,m range, 180\textdegree{} front-facing
field of view) with the polar representation of the goal direction
($d$: Euclidean distance, $\cos\theta, \sin\theta$: relative bearing)
and the previous action. The action $a = [v, \omega]$ specifies the
linear and angular velocity of a differential-drive robot. The reward
function follows a sparse formulation augmented with a local shaping
term:
\begin{equation}
r = \begin{cases}
+100 & \text{goal reached} \\
-100 & \text{collision} \\
v - \frac{|\omega|}{2} - \frac{1}{2}\max(0,\, 1.35 - d_{min}) & \text{otherwise}
\end{cases}
\end{equation}
where $d_{min}$ is the minimum LiDAR return, encouraging forward
progress while penalizing sharp turns and close proximity to obstacles.

\subsection{Baseline Architectures}

\textbf{CNNTD3.} The actor network processes the LiDAR scan through
three 1D convolutional layers
(Conv1D$(1{\to}4, k{=}8, s{=}4)$ $\to$
Conv1D$(4{\to}8, k{=}8, s{=}4)$ $\to$
Conv1D$(8{\to}4, k{=}4, s{=}2)$), producing a 16-dimensional feature
vector. This is concatenated with a goal embedding
($3{\to}10$, linear) and a previous-action embedding ($2{\to}10$,
linear), yielding a 36-dimensional joint representation that is passed
through two fully connected layers (400, 300 units) before a
\texttt{tanh}-bounded 2-dimensional action output. The critic follows
the standard twin-Q architecture of TD3~\cite{fujimoto2018td3}, sharing
the same LiDAR feature extractor.

\textbf{RCPG.} As a memory-augmented baseline, we replace the
convolutional feature extractor with a GRU operating over a 10-step
history of raw states (hidden size 256), testing whether temporal
context improves navigation in structured environments.

\subsection{Curriculum Learning with Exploration Reward}

Standard random-obstacle training never exposes the agent to concave
trap geometries (e.g., U-shaped enclosures), where the goal-direction
signal points directly through a wall. We address this with two
complementary mechanisms.

\textbf{Curriculum scheduling.} Starting from epoch 10, each episode
is drawn from a structured hard-scenario pool
(U-trap, U-shape, U-shape-hard) with probability
$p_{hard} = \min(0.5,\ \text{epoch}/E_{max})$, linearly increasing the
exposure to difficult geometries while preserving sufficient standard-
environment training to avoid catastrophic forgetting of baseline
navigation behavior.

\textbf{Exploration reward.} We augment the reward with two terms that
penalize spatial stagnation, computed from a sliding window of the
agent's last 50 positions $\mathcal{P}$:
\begin{equation}
r_{explore} =
\begin{cases}
+0.3 & \min_{p \in \mathcal{P}} \|x_t - p\| > 0.5\,\text{m} \\
-0.2 & \min_{p \in \mathcal{P}} \|x_t - p\| < 0.1\,\text{m} \\
0 & \text{otherwise}
\end{cases}
\end{equation}
\begin{equation}
r_{stall} =
\begin{cases}
-0.5 & \text{stall\_count} > 15 \\
-0.2 & \text{stall\_count} > 8 \\
0 & \text{otherwise}
\end{cases}
\end{equation}
where \texttt{stall\_count} tracks consecutive steps with displacement
below 2\,cm. The combined reward $r + r_{explore} + r_{stall}$ is used
only for the replay buffer target; evaluation always uses the
standard-environment-only reward to ensure fair comparison across
methods.

We empirically found exploration reward magnitude to be a sensitive
hyperparameter: a moderate setting (0.3 / $-$0.2, denoted
\textbf{CNNTD3-Improved}) successfully escapes U-trap scenarios but at
a 2-point cost to standard-environment success rate, while weaker
settings (0.1 / $-$0.05, denoted \textbf{CNNTD3-v3}) and an
intermediate setting (0.15 / $-$0.1) fail to consistently escape traps
despite longer training (100 and 80 epochs respectively), suggesting a
threshold effect rather than a monotonic trade-off. This finding is
analyzed further in Section~\ref{sec:results} as an ablation.
